% Chapitre 2 : Méthodologie et bonnes pratiques

\section{Méthodologie et bonnes pratiques}

\subsection{Planification d'un projet d'annotation}

La réussite d'un projet d'annotation nécessite une planification rigoureuse :

\subsubsection{Définition des objectifs}
\begin{itemize}
    \item Clarifier les objectifs du projet d'annotation
    \item Identifier les cas d'usage finaux des données annotées
    \item Définir les métriques de qualité et de succès
\end{itemize}

\subsubsection{Conception du schéma d'annotation}
\begin{itemize}
    \item Élaborer un schéma d'annotation adapté aux objectifs
    \item Définir les guidelines précises pour les annotateurs
    \item Prévoir des tests pilotes pour valider le schéma
\end{itemize}

\subsection{Assurance qualité}

\subsubsection{Accord inter-annotateurs}
L'accord inter-annotateurs mesure la cohérence des annotations entre différents annotateurs :

\begin{itemize}
    \item \textbf{Coefficient de Kappa} : Mesure standard de l'accord
    \item \textbf{F1-score} : Évaluation de la précision et du rappel
    \item \textbf{Seuils d'acceptabilité} : Définir des seuils minimums d'accord
\end{itemize}

\subsubsection{Processus de validation}
\begin{itemize}
    \item Sessions de formation des annotateurs
    \item Annotations de référence (gold standard)
    \item Révisions périodiques et corrections
    \item Contrôle qualité en continu
\end{itemize}

\subsection{Gestion des annotateurs}

\subsubsection{Formation et accompagnement}
\begin{itemize}
    \item Formation initiale sur l'outil et les guidelines
    \item Sessions de calibration régulières
    \item Support technique et méthodologique
\end{itemize}

\subsubsection{Suivi de la performance}
\begin{itemize}
    \item Métriques de productivité
    \item Indicateurs de qualité individuels
    \item Feedback constructif et amélioration continue
\end{itemize}

\subsection{Bonnes pratiques d'annotation}

\subsubsection{Cohérence}
\begin{itemize}
    \item Appliquer les guidelines de manière uniforme
    \item Résoudre les ambiguïtés par consensus
    \item Documenter les cas particuliers
\end{itemize}

\subsubsection{Efficacité}
\begin{itemize}
    \item Utiliser les raccourcis clavier
    \item Optimiser l'organisation du workspace
    \item Planifier les sessions d'annotation
\end{itemize}

\newpage